\documentclass[11pt,letterpaper]{article}
\usepackage[margin=1in,headheight=14pt]{geometry}
\usepackage[T1]{fontenc}
\usepackage{lmodern}
\usepackage{amsmath,amssymb,amsthm,mathtools}
\usepackage{microtype,booktabs,xcolor,fancyhdr}
\definecolor{ink}{HTML}{19354A}
\usepackage[colorlinks=true,linkcolor=ink,urlcolor=ink,citecolor=ink]{hyperref}
\hypersetup{pdftitle={Matching foundations and exact small-host studies},pdfauthor={Akhil Ravuri}}
\newtheorem{theorem}{Theorem}
\newtheorem{corollary}[theorem]{Corollary}
\pagestyle{fancy}\fancyhf{}
\fancyhead[L]{\small\textsc{Matching foundations}}
\fancyhead[R]{\small Akhil Ravuri}\fancyfoot[C]{\thepage}
\renewcommand{\headrulewidth}{0.3pt}
\setlength{\parskip}{0.2em}\setlength{\emergencystretch}{1em}
\title{\vspace{-1.5em}\LARGE\bfseries Matching foundations\\and exact small-host studies}
\author{Akhil Ravuri}\date{}
\begin{document}
\maketitle\vspace{-1.5em}
\begin{abstract}
We record the standard regular-bipartite corollary of Hall's theorem, its
finite algorithm checks, and small rainbow-perfect-matching searches. The
$K_{3,3}$ experiment classifies 87 extremal fixed-edge partitions into four
symmetry classes. General matching thresholds quoted here are known results;
the classifications are finite computations, without a novelty claim.
For the first three ladders, we prove the maximum color counts and compute
all extremizers and their exact random-coloring probabilities by two methods.
\end{abstract}

\section{Hall's theorem and a regular-graph corollary}
Let $G=(L,R,E)$ be a finite bipartite graph. For $S\subseteq L$, let $N(S)$
be the set of its neighbors. Hall's marriage theorem states that a matching
saturating $L$ exists if and only if $|N(S)|\ge|S|$ for every $S\subseteq L$.
The matching-deficiency formulation in~\cite[Lemma 2.1]{LTJ} is equivalent.
We use this known theorem, rather than prove it here.

\begin{corollary}
Every finite $d$-regular bipartite graph with $d\ge1$ has a perfect matching.
\end{corollary}
\begin{proof}
Counting edges from the two sides gives $d|L|=|E|=d|R|$, hence $|L|=|R|$.
For any $S\subseteq L$, exactly $d|S|$ edges run from $S$ into $N(S)$, while
at most $d|N(S)|$ edges can enter those neighbors. Since $d>0$,
$d|S|\le d|N(S)|$ implies $|S|\le|N(S)|$. Hall's theorem gives a matching
saturating $L$; equal part sizes make it perfect.
\end{proof}

The generator chooses a $d$-element neighbor set for each left vertex and
retains exactly those choices giving every right vertex degree $d$. Thus each
labeled graph appears once. For $1\le n\le5$ and $1\le d\le n$, the row
counts, ordered by increasing $d$, are
\begin{center}
\begin{tabular}{crr}
\toprule
Part size $n$ & Counts for $d=1,\ldots,n$ & Total\\
\midrule
1 & $1$ & 1\\
2 & $2,1$ & 3\\
3 & $6,6,1$ & 13\\
4 & $24,90,24,1$ & 139\\
5 & $120,2040,2040,120,1$ & 4321\\
\midrule
 & & 4477\\
\bottomrule
\end{tabular}
\end{center}
An augmenting-path routine finds a perfect matching in all 4,477 cases.
Independent full-mask generation and partial-injection matching oracles check
the implementation through part size three. The computation tests the code;
the corollary's proof does not depend on it.

\section{Rainbow conventions and known thresholds}
For a fixed graph $G$ with a perfect matching of at least two edges, define
\[
 A(G)=\max\{|c(E(G))|:\text{$c$ has no rainbow perfect matching}\}.
\]
Here colorings need not be proper. An \emph{exact} $r$-coloring uses all $r$
colors. In the feasible cases considered here, the minimum exact-color
forcing threshold is $A(G)+1$: merging classes preserves avoidance, so an
extremal avoiding coloring supplies every smaller positive color count.

Li, Tu, and Jin~\cite[Theorem 2.5]{LTJ} prove that, for $m\ge n\ge k\ge3$,
\[
 \operatorname{rb}(K_{m,n},kK_2)=m(k-2)+2.
\]
Their $\operatorname{rb}$ is a minimum forcing number. Consequently
$A(K_{n,n})=n(n-2)+1$ for $n\ge3$. The case $n=2$ uses their separate
Theorem 2.3, giving $A(K_{2,2})=2$. In particular, the repository's values
$2$ and $4$ are reproductions of known results, not conjectures.

For the lower-bound construction when $n\ge3$, assign distinct colors to all
edges incident to $n-2$ left vertices and one fresh common color to all
remaining edges. This uses $n(n-2)+1$ colors, and every perfect matching uses
two edges of the common color. The cited theorem supplies the upper bound.

\section{Four symmetry classes on \texorpdfstring{$K_{3,3}$}{K3,3}}
The search examines all $B_9=21{,}147$ partitions of nine row-major edge
positions. Exactly 87 avoiding partitions use the maximum of four colors.
These counts keep edge positions fixed and ignore only color names.

For the symmetry quotient, independently permute rows and columns and allow
matrix transpose, giving the full automorphism group of order
$2(3!)^2=72$. Normalize color names by first appearance after each action;
the lexicographically least normalized array is the representative.
\begin{center}\small
\begin{tabular}{clrr}
\toprule
Class & Representative rows & Orbit & Stabilizer\\
\midrule
1 & \texttt{000 / 000 / 123} & 6 & 12\\
2 & \texttt{001 / 110 / 230} & 36 & 2\\
3 & \texttt{001 / 120 / 230} & 36 & 2\\
4 & \texttt{001 / 230 / 320} & 9 & 8\\
\bottomrule
\end{tabular}
\end{center}
The orbit sizes sum to 87; each orbit-stabilizer product equals 72. The last
two classes have the same color-class size multiset $\{1,2,2,4\}$, so that
multiset alone does not classify the partitions. Stabilizers fix the partition
and may permute color names. Every transformed partition is checked to remain
inside the enumerated extremal population, and every canonical bucket is
compared with its explicit orbit. This does not certify absence from prior
literature or classify larger hosts.

\newpage
\section{Sparse-host comparison and exact coverage}
The separate sparse-host search gives the following fixed-edge counts.
\begin{center}\small
\begin{tabular}{lrrrrr}
\toprule
Host & Edges & Matchings & Partitions covered & $A(G)$ & Extremizers\\
\midrule
$C_4$ & 4 & 2 & 15 & 2 & 1\\
$C_6$ & 6 & 2 & 203 & 4 & 9\\
$C_8$ & 8 & 2 & 4140 & 6 & 36\\
$C_4\sqcup C_4$ & 8 & 4 & 4140 & 6 & 2\\
$K_{3,3}$ & 9 & 6 & 21147 & 4 & 87\\
$Q_3$ & 12 & 9 & 4213597 & 7 & 174\\
\bottomrule
\end{tabular}
\end{center}
The cycle and disjoint-square companion paper gives elementary extremal
proofs. The cube companion paper proves its maximum and records the nine
computed symmetry classes; 174 counts fixed-edge partitions, not classes.

The solver enumerates restricted-growth strings. When a completed perfect
matching is rainbow in a prefix, every continuation remains forbidden.
If $T(r,b)$ counts completions of $r$ positions with $b$ existing classes, then
\[
 T(0,b)=1,\qquad T(r,b)=bT(r-1,b)+T(r-1,b+1).
\]
Rejected subtrees are counted by this recurrence. Surviving leaves plus
disjoint rejected subtrees equal $B_m=T(m-1,1)$ for $m$ edges. Thus coverage
does not mean visiting every leaf individually. An independent unpruned
partition oracle checks every balanced bipartite host of part size two or
three with a perfect matching; edge-subset oracles check matching constraints.

\subsection*{Reproduction and records}
From the repository root, use Python 3.11+ and its standard library:
\begin{quote}\small\ttfamily
python -m experiments.reproduce\par
python -m experiments.sparse\_rainbow\par
python -m experiments.classify\_k33
\end{quote}
Reference datasets are \path{results/data/small_cases.json} and
\path{results/data/sparse_rainbow.json}. The
\href{https://github.com/Akhil-rav/discrete-math-lab/blob/main/results/verification_index.md}{verification index}
links execution records, independent checks, and precise ranges. These
finite searches and elementary reproductions make no novelty claim.

\newpage
\section{Three small ladders}
Let $L_k=P_k\mathbin{\square}K_2$, with two rows and $k$ columns. Number columns
from zero and write $r_i$ for a rung, $t_i$ for the top edge across gap $i$,
and $b_i$ for the bottom edge across that gap. Edge positions are ordered
\[
 r_0,\ldots,r_{k-1},\quad
 t_0,\ldots,t_{k-2},\quad b_0,\ldots,b_{k-2}.
\]
There are $m=3k-2$ edges. A perfect matching takes either the first rung
or both rail edges across the first gap; the remaining graph is a shorter
ladder. Thus the usual domino-tiling recurrence is
$p_k=p_{k-1}+p_{k-2}$, with $p_0=p_1=1$. The first three hosts have
2, 3, and 5 perfect matchings. In edge-index notation, the larger lists are
\[
\begin{array}{c|l}
 L_3 & 012,\ 046,\ 235\\
 L_4 & 0123,\ 0169,\ 0358,\ 2347,\ 4679 .
\end{array}
\]
Here $046$ denotes the edge set $\{0,4,6\}$, not a coloring.

\begin{theorem}
$A(L_2)=2$, $A(L_3)=4$, and $A(L_4)=7$.
\end{theorem}
\begin{proof}
For $L_2=C_4$, each of its two alternating matchings must be monochromatic;
using a different color on each gives the maximum of two.

In $L_3$, repeat a distinct color on each of the disjoint pairs
$\{r_0,r_1\}$, $\{t_0,b_0\}$, $\{t_1,b_1\}$, and give $r_2$ a fourth color.
Every perfect matching contains one of these pairs. In $L_4$, repeat a
distinct color on each of $\{r_0,r_3\}$, $\{t_0,b_0\}$, $\{t_2,b_2\}$,
and give the four remaining edges fresh colors. The displayed matching list
shows that every matching again contains a repeated pair. These constructions
give four and seven colors, respectively.

For the upper bounds, an avoiding coloring with at least $m-2$ colors could
be merged down to exactly $m-2$ colors while preserving avoidance. Its
deficit $\sum_B(|B|-1)=2$ forces either one triple or two doubletons,
with every other class a singleton.

For $L_3$, any fixed pair belongs to at most one perfect matching, so two
doubletons obstruct at most two of the three matchings. Each edge belongs
to at most two matchings; only $r_0,r_2$ attain two. Hence any triple has
at most $2+2+1=5$ edge--matching incidences. Obstructing all three
matchings would require at least six incidences.

For $L_4$, any pair belongs to at most two matchings, so two doubletons
obstruct at most four of the five. The two end rungs have incidence three;
every other edge has incidence at most two. A triple therefore has at
most $3+3+2=8$ incidences, fewer than the ten required to obstruct all five.
All these incidence bounds can be read directly from the complete lists.
Thus $A(L_3)\le7-3=4$ and $A(L_4)\le10-3=7$.
\end{proof}
The bounds are elementary proofs, independent of the coloring enumerator.
They do not assert a formula for arbitrary ladder length.

\newpage
\section{Complete finite ladder classifications}
All $B_4=15$, $B_7=877$, and $B_{10}=115{,}975$ edge partitions were visited
without rainbow pruning. Their avoiding histograms agree with the separate
pruned solver. There are 1, 27, and 25 extremizers, respectively.
Color names are normalized by first appearance. The table lists the complete
symmetry quotient; strings use the edge order fixed above.
\begin{center}\small
\begin{tabular}{clrr}
\toprule
Host & Representative & Orbit size & Stabilizer size\\
\midrule
$L_2$ & \texttt{0011} & 1 & 8\\
\midrule
$L_3$ & \texttt{0011023} & 4 & 1\\
 & \texttt{0011230} & 4 & 1\\
 & \texttt{0011232} & 4 & 1\\
 & \texttt{0012023} & 4 & 1\\
 & \texttt{0012323} & 2 & 2\\
 & \texttt{0100023} & 2 & 2\\
 & \texttt{0100230} & 2 & 2\\
 & \texttt{0100232} & 4 & 1\\
 & \texttt{0102323} & 1 & 4\\
\midrule
$L_4$ & \texttt{0010234256} & 2 & 2\\
 & \texttt{0012304356} & 4 & 1\\
 & \texttt{0012324356} & 4 & 1\\
 & \texttt{0012345346} & 2 & 2\\
 & \texttt{0120030456} & 2 & 2\\
 & \texttt{0120034560} & 2 & 2\\
 & \texttt{0120034564} & 4 & 1\\
 & \texttt{0120234564} & 4 & 1\\
 & \texttt{0120345365} & 1 & 4\\
\bottomrule
\end{tabular}
\end{center}
The full automorphism groups have orders 8, 4, and 4. For $L_3,L_4$ the
actions are generated by exchanging the two rails and reversing the column
order; for $L_2$ all square symmetries are used. Enumeration of all vertex
bijections independently recovers these groups. Every orbit is checked to
lie in the enumerated extremal population, different recorded orbits are
disjoint, and their union equals that population. Stabilizers fix partitions,
allowing color names to permute.

Both larger hosts have color deficit three, so the only nonsingleton
class-size patterns are a quadruple, a triple and a pair, or three pairs.
Their fixed-edge counts, in this order, are $(4,16,7)$ for $L_3$ and
$(4,10,11)$ for $L_4$. Nine symmetry classes in both cases do not imply
the same fixed-edge population or a general classification theorem.
All partitions, class representatives, and orbit sizes are saved in
\path{results/data/ladder_matchings.json}.

\subsection*{Reproduction and source scope}
From the repository root, run
\begin{quote}\small\ttfamily
python -m experiments.ladder\_matchings
\end{quote}
Use \texttt{--check} with the dataset path above to compare all saved data.
The run and precise test ranges are in \path{results/ladder_verification.md};
the limited source comparison is in \path{notes/literature/ladders.md}.

\newpage
\section{Exact ladder probabilities and checks}
Independently color every edge uniformly from $q\ge1$ available labels; using
every label is not required. If $a_j$ counts avoiding partitions with $j$
classes and $(q)_j=q(q-1)\cdots(q-j+1)$, then
\[
 P_k(q)=1-\frac{\sum_j a_j(q)_j}{q^{3k-2}} .
\]
Each partition has exactly $(q)_j$ assignments of distinct labels to its
classes. The computed histograms, including every positive entry, are
\[
\begin{array}{c|rrrrrrr}
 & a_1&a_2&a_3&a_4&a_5&a_6&a_7\\ \hline
L_2&1&1&0&0&0&0&0\\
L_3&1&63&108&27&0&0&0\\
L_4&1&511&9330&17629&7710&948&25 .
\end{array}
\]
Consequently the rainbow-coloring count polynomials $U_k(q)=q^{3k-2}P_k(q)$
are
\begin{align*}
U_2(q)&=q^4-q^2,\\
U_3(q)&=q^7-27q^4+54q^3-36q^2+8q,\\
U_4(q)&=q^{10}-25q^7-423q^6+2135q^5\\
 &\qquad{}-2734q^4-706q^3+3408q^2-1656q .
\end{align*}
The separate matching-event method forms a conflict graph on edge positions
for each nonempty subfamily of perfect matchings, counts its proper colorings,
and applies inclusion--exclusion. It handles 3, 7, and 31 event subfamilies
on these hosts. The resulting polynomials agree coefficient by coefficient
with the partition-derived expressions, and both methods give identical
exact probabilities for every $1\le q\le24$. The probability companion paper
proves this general counting method.

For four available labels,
\[
 P_2(4)=\frac{15}{16},\qquad
 P_3(4)=\frac{387}{512},\qquad
 P_4(4)=\frac{12357}{32768}.
\]
For both larger ladders the avoidance probability is of order $q^{-3}$:
$q^3(1-P_3(q))\to27$ and $q^3(1-P_4(q))\to25$ as $q\to\infty$.
Indeed the highest nonzero histogram term has degree $m-3$, and each falling
factorial is monic. These are limits derived from the finite polynomials.

\begin{thebibliography}{1}
\bibitem{LTJ} X. Li, J. Tu, and Z. Jin,
\emph{Bipartite rainbow numbers of matchings}, Discrete Math.
\textbf{309} (2009), 2575--2578.
\href{https://arxiv.org/abs/math/0610910v1}{arXiv:math/0610910v1}.
Definitions and Lemma 2.1, Theorems 2.3 and 2.5 checked in this preprint;
see the repository's source records for notation translation.
\end{thebibliography}
\end{document}
